
\begin{proposition}[Every compact elliptic surface satisfies the intrinsic rank-one field condition]\label{cdp:st:C.2.UN.26}
Let $q:X\to B$ be a proper holomorphic map from a
smooth compact connected complex surface to a smooth
compact curve, with connected elliptic general fibre.
There are commuting meromorphic vector fields $V_1,V_2$
which are linearly independent over $\mathbb C$, have
generic rank one, and whose every constant linear
combination is semicomplete on its full holomorphic
locus. Their ratio has an actual graph resolution and
Stein factor birationally equal to $q$.

If a finite cyclic group $G$ acts holomorphically on
$X$ and preserves $q$, the two-dimensional space can
be chosen $G$-stable. Thus the rank-one field-space
condition of \ref{cdp:st:C.2.UN.16}, even with this equivariance,
is satisfied by every specified compact elliptic
surface and cyclic fibration action. Any additional
restriction on its occurrence as a degeneration
component must use further data of that degeneration.
\end{proposition}

\begin{proof}
The relative canonical line is the actual invertible
sheaf
\[
 \omega_{X/B}=K_X\otimes q^*K_B^{-1},
 \qquad \mathcal E=q_*\omega_{X/B}.
                                                               \tag{C.FL.1}\label{cdp:eq:C.FL.1}
\]
Proper coherence makes $\mathcal E$ coherent. It is
torsion-free: a section of a line bundle annihilated
by a nonzero base function vanishes over the dense
complement of its zero fibre and hence vanishes
identically. On a smooth elliptic fibre the canonical
space has dimension one, so $\mathcal E$ has rank one.
It is therefore a line bundle on $B$.

Choose a nonzero meromorphic section $\sigma$ of
$\mathcal E$. This is an actual global choice:
Riemann--Roch on $B$ gives a nonzero holomorphic
section of $\mathcal E(D)$ for an effective divisor
$D$ of sufficiently large degree, which is the
required meromorphic section of $\mathcal E$.
Evaluation gives a nonzero meromorphic section
\[
 a:\mathcal O_X\dashrightarrow\omega_{X/B}
       \quad\hbox{from}\quad q^*\mathcal E\longrightarrow
                                      \omega_{X/B}.
                                                               \tag{C.FL.2}\label{cdp:eq:C.FL.2}
\]
Here the dashed notation records poles of the section,
not a replacement of the displayed sheaf evaluation.
On the smooth part of $q$, the canonical differential
identification $\omega_{X/B}=\Omega^1_{X/B}$ makes
$a$ an actual relative one-form.

Its dual defines a global meromorphic vertical field
$V$. The coordinate comparison at critical fibres
can be specified as follows. On a coordinate disc
of $B$ with coordinate $z$, set
$\Omega_z=a\otimes q^*(dz)$, a meromorphic section
of $K_X$. Define $V$ by
\[
                   \iota_V\Omega_z=-q^*(dz).
                                                               \tag{C.FL.3}\label{cdp:eq:C.FL.3}
\]
Contraction by a nonzero meromorphic two-form is an
isomorphism between the meromorphic tangent and
cotangent bundles, so this has a unique meromorphic
solution. Under $z'=\phi(z)$ both sides are multiplied
by $q^*\phi'(z)$; the solution is unchanged and
therefore glues. The alternating identity in
\eqref{cdp:eq:C.FL.3} gives $dq(V)=0$. On a regular fibre its
relative time form is $a$, so $a(V)=1$ with the
usual convention $\Omega_z=q^*(dz)\wedge a$.

Remove finitely many base values so that $q$ is
smooth and the evaluated section is holomorphic
and nonzero on every remaining fibre. Such a
removal is possible by generic base change for
$q_*\omega_{X/B}$ and the finite zero and pole
divisors of $\sigma$. A nonzero canonical form
on a compact elliptic curve is nowhere zero.
Thus $V$ restricts on each of these fibres to
a nonzero holomorphic translation field and is
complete there.

Choose any nonconstant meromorphic function
$h$ on $B$, and put
\[
              V_1=V,\qquad V_2=(q^*h)V.
                                                               \tag{C.FL.4}\label{cdp:eq:C.FL.4}
\]
Their bracket vanishes because $V(q^*h)=0$.
They are independent over $\mathbb C$ since $h$
is nonconstant, and have generic rank one by
the formula. For $c,d\in\mathbb C$, the field
$W=(c+dq^*h)V$ is complete on the dense union
of the preceding compact smooth fibres away
from the poles of $h$. Its coefficient is
constant on each such fibre; zero coefficients
give constant complete solutions. Apply
\cite[Theorem~3.3]{Guillot2024} to the complex
manifold obtained by deleting the actual polar
divisor of $W$. The univalence locus is closed
there and contains this dense union. It is
therefore the whole holomorphic locus, proving
the asserted semicompleteness for every $W$.

The ratio is $q^*h$, so the resulting map is
the actual composite $X\xrightarrow q B
\xrightarrow h\mathbb P^1$. The second arrow
is finite. Its Stein factor is $q$ itself:
$q_*\mathcal O_X=\mathcal O_B$, and $B$ is
the normalization in the field of functions
constant on the connected general fibres.
This proves the precise first-integral comparison.

For the equivariant assertion, the finite orbit
of $V$ spans a nonzero finite-dimensional
representation of $G$ by meromorphic vertical
fields. A cyclic complex representation has
an eigenvector $V'\ne0$. On a general smooth
fibre each orbit field is holomorphic and
vertical, and therefore a translation field.
The same is true of $V'$; after removing
finitely many further values it is nonzero
on every remaining fibre. Choose a nonconstant
meromorphic function $h$ on the quotient
curve $B/G$, and pull it back to $B$.
Then the space
\[
                  \mathbb C V'+\mathbb C(q^*h)V'
                                                               \tag{C.FL.5}\label{cdp:eq:C.FL.5}
\]
is $G$-stable, with the same eigencharacter
on its two displayed generators. Its members
are semicomplete by the dense compact-fibre
and closedness argument just proved, and its
ratio has the same connected-fibre Stein
factor $q$. This proves the full statement.
\end{proof}


